\documentclass[11pt]{amsart}
\usepackage[T1]{fontenc}
\usepackage{amsmath,amssymb,amsthm,mathtools}
\usepackage{graphicx}
\usepackage[colorlinks=true,linkcolor=blue,citecolor=blue,urlcolor=blue]{hyperref}

\newtheorem{theorem}{Theorem}
\newtheorem{lemma}{Lemma}
\newtheorem{corollary}{Corollary}
\newcommand{\M}{\mathcal M}
\newcommand{\Tr}{\operatorname{Tr}}
\newcommand{\gammatwo}{\gamma_2}
\newcommand{\sa}{\mathrm{sa}}

\begin{document}

\title{Entangled Werner States Undetected by Every Entanglement Tester}
\author{Run fa\_banach\_001}
\date{August 11, 2026}
\maketitle
\pagestyle{plain}

\section*{Status}

This is a candidate counterexample, subject to human review, to the main open
question of Jivulescu, Lancien, and Nechita \cite{JLN}.  Without the index
permutation allowed in their positive theorem, the family of entanglement
testers is not complete.  In fact, a whole interval of entangled Werner
states is invisible to every possible tailored pair of testers.

\section*{Source Question}

For a multipartite state $\rho$, the source applies contractions
$\mathcal E_i:S_1^{d_i}\to\ell_2^{n_i}$ and tests whether
\[
 \| (\mathcal E_1\otimes\cdots\otimes\mathcal E_m)(\rho)\|_
 {\ell_2^{n_1}\otimes_\pi\cdots\otimes_\pi\ell_2^{n_m}}>1.
\]
It asks whether every entangled state is detected in this way, without first
permuting indices.  The source notes that in the bipartite case this asks
whether every entangled state has a norm-one witness factoring through
$\ell_2$ \cite[Section 13, p.~35]{JLN}.

\begin{center}
\includegraphics[width=\textwidth]{figures/open_problem_crop.png}
\end{center}

\section*{Proof Intuition}

The source already identifies the right geometric object: tester witnesses
are maps $S_1^d\to S_\infty^d$ with Hilbert-factorization norm at most one.
Werner states are invariant under $U\otimes U$.  Averaging any such witness
over this symmetry preserves both its value on the state and its
factorization bound, and leaves only a two-parameter witness $aI+bF$.

The decisive calculation is exact.  For the map represented by $aI+bF$,
\[
 \gammatwo(T_{a,b})
 =\max\left\{|b|,
 \left|a+\frac bd\right|+\left(1-\frac1d\right)|b|\right\}.
\]
The upper bound splits a matrix into its scalar and traceless parts and sends
the latter through Hilbert--Schmidt space.  For the lower bound, an
off-diagonal matrix unit sees $|b|$, while restriction to diagonal matrices
produces $aJ+bI$; its nuclear norm divided by $d$ gives the other term.
The remaining two-variable optimization shows that no tailored tester beats
ordinary realignment on a Werner state.

\section*{Factorable Witnesses}
\par\medskip
Use the complex bilinear trace duality between $S_1^d$ and $S_\infty^d$.
For $W\in\M_d\otimes\M_d$, let $T_W:S_1^d\to S_\infty^d$ be determined by
\[
 \Tr\bigl(T_W(X)Y\bigr)=\Tr\bigl(W(X\otimes Y)\bigr).
\]
Changing the tensor-to-map convention can insert a transpose, which is an
isometry on both Schatten classes and does not affect any norm below.  Put
\[
 \gammatwo(T)=\inf\{\|A\|\,\|B\|:
 T=BA,\ A:S_1^d\to H,\ B:H\to S_\infty^d,
 \ H\text{ Hilbert}\}.
\]

For a state $\rho$, define its all-tester envelope
\[
 \alpha(\rho)=\sup_{\mathcal E,\mathcal F}
 \|(\mathcal E\otimes\mathcal F)(\rho)\|_\pi,
\]
where the supremum is over contractions from the two local $S_1$ spaces to
arbitrary finite-dimensional Hilbert spaces.  Requiring the maps to have norm
exactly one gives the same detection question.

\begin{lemma}[Tester duality]\label{lem:duality}
For every bipartite state $\rho$,
\[
 \alpha(\rho)\leq
 \sup_{\gammatwo(T_W)\leq1}|\Tr(W\rho)|.
\]
\end{lemma}

\begin{proof}
Fix contractions
\[
 \mathcal E:S_1^d\to H,\qquad \mathcal F:S_1^d\to K.
\]
By projective/injective duality on $H\otimes K$,
each functional used to norm $(\mathcal E\otimes\mathcal F)(\rho)$ is a
bounded bilinear form $v$ on $H\times K$ of norm at most one.  By the Riesz
representation theorem, $v$ is represented by a contraction between the two
Hilbert spaces.  Hence
\[
 (X,Y)\longmapsto v(\mathcal E X,\mathcal F Y)
\]
has associated map $S_1^d\to S_\infty^d$ factoring through a Hilbert space
with product of norms at most one.  Its representing tensor is some $W$ with
$\gammatwo(T_W)\leq1$.  Taking the successive suprema proves the claim.
\end{proof}

\section*{The Exact Invariant Factorization Norm}

Let $F(x\otimes y)=y\otimes x$ be the flip and put
$W_{a,b}=aI+bF$.  Then, up to the harmless transpose mentioned above,
\[
 T_{a,b}(X)=a\Tr(X)I+bX.
\]

\begin{lemma}\label{lem:gamma}
For every $a,b\in\mathbb C$ and $d\geq2$,
\begin{equation}\label{eq:gamma}
 \gammatwo(T_{a,b})=
 \max\left\{|b|,
 \left|a+\frac bd\right|+\left(1-\frac1d\right)|b|\right\}.
\end{equation}
\end{lemma}

\begin{proof}
Write
\[
 A=a+\frac bd,\qquad c=1-\frac1d,
 \qquad X_0=X-\frac{\Tr X}{d}I.
\]
For $r,s>0$, factor $T_{a,b}=G_{r,s}E_{r,s}$ through
$\mathbb C\oplus_2 S_2^{d,0}$ by
\[
 E_{r,s}(X)=(r\Tr X,sX_0),\qquad
 G_{r,s}(z,Z)=\frac Ar zI+\frac bs Z.
\]
The extreme points of the complex $S_1$ unit ball are rank-one matrices
$uv^*$ with $\|u\|=\|v\|=1$.  If $t=\Tr(uv^*)$, then
$\|(uv^*)_0\|_2^2=1-|t|^2/d$.  It follows, and by duality the same calculation
applies to $G_{r,s}$, that
\begin{align}
 \|E_{r,s}\|^2&=\max\{s^2,r^2+cs^2\},\label{eq:E}\\
 \|G_{r,s}\|^2&=\max\left\{\frac{|b|^2}{s^2},
 \frac{|A|^2}{r^2}+c\frac{|b|^2}{s^2}\right\}.
 \label{eq:G}
\end{align}
If $b=0$, the scalar factorization gives $\gammatwo(T_{a,0})\leq|a|$.
Assume $b\ne0$ and set $B=|b|$, $C=|A|$, and $s=1$.  If $C\leq B/d$,
choose $r^2$ in the nonempty interval
\[
 d(C/B)^2\leq r^2\leq1/d.
\]
Equations \eqref{eq:E}--\eqref{eq:G} give
$\|E_{r,1}\|\|G_{r,1}\|=B$.  If $C\geq B/d$, choose
$r^2=C/B$; the product is $C+cB$.  This proves the upper bound in
\eqref{eq:gamma}.

For the reverse bound, $\gammatwo(T)\geq\|T\|$.  Applying $T_{a,b}$ to an
off-diagonal matrix unit gives $\gammatwo(T_{a,b})\geq|b|$.  Next restrict the
domain to diagonal matrices and project the range onto its diagonal.  These
maps have norm one, so the resulting operator
\[
 M=aJ+bI:\ell_1^d\longrightarrow\ell_\infty^d
\]
has $\gammatwo(M)\leq\gammatwo(T_{a,b})$.  If
$M_{ij}=\langle x_j,y_i\rangle$ is any Hilbert factorization, then
\[
 \|M\|_*\leq
 \left(\sum_j\|x_j\|^2\right)^{1/2}
 \left(\sum_i\|y_i\|^2\right)^{1/2}
 \leq d\,\|A_0\|\,\|B_0\|,
\]
where $A_0:\ell_1^d\to H$ and $B_0:H\to\ell_\infty^d$ are the two
factors.  Thus $\gammatwo(M)\geq\|M\|_*/d$.  The eigenvalues (and singular
values in modulus) of $aJ+bI$ are $dA$ once and $b$ with multiplicity
$d-1$.  Consequently
\[
 \gammatwo(T_{a,b})\geq\frac{\|M\|_*}{d}
 =|A|+c|b|.
\]
Together with the off-diagonal bound, this proves equality.
\end{proof}

\section*{Werner States and the Negative Answer}

For $0\leq\mu\leq1$, the Werner state \cite{Werner,JLN} is
\begin{equation}\label{eq:werner}
 \sigma_\mu=
 \mu\frac{I+F}{d(d+1)}+(1-\mu)\frac{I-F}{d(d-1)}.
\end{equation}
It is positive, has trace one, is invariant under $U\otimes U$, and satisfies
\begin{equation}\label{eq:f}
 f_\mu:=\Tr(F\sigma_\mu)=2\mu-1.
\end{equation}

\begin{theorem}[Exact all-tester envelope]\label{thm:envelope}
For every $d\geq2$ and $0\leq\mu\leq1$,
\begin{equation}\label{eq:envelope}
 \alpha(\sigma_\mu)=1+2\left(\frac1d-\mu\right)_+.
\end{equation}
Hence the realignment tester is optimal among all tailored pairs of
entanglement testers on every Werner state.
\end{theorem}

\begin{proof}
The $\gammatwo$ unit ball is convex and invariant under local unitary
conjugations.  Given a witness $W$ in that ball, average it over
$W\mapsto(U\otimes U)W(U^*\otimes U^*)$.  Its value on $\sigma_\mu$ is
unchanged, and the average remains in the ball.  Taking a Hermitian part also
does not increase $\gammatwo$.  By the elementary two-fold Schur--Weyl
commutant, the averaged Hermitian witness has the form $aI+bF$ with
$a,b\in\mathbb R$.

Put $A=a+b/d$ and $c=1-1/d$.  Lemma \ref{lem:gamma} says that the feasible
region is
\[
 |b|\leq1,\qquad |A|+c|b|\leq1.
\]
Using \eqref{eq:f}, the objective is
\[
 \Tr((aI+bF)\sigma_\mu)=A+b(f_\mu-1/d).
\]
For fixed $b$ it is maximal at $A=1-c|b|$.  Maximizing over
$-1\leq b\leq1$ gives
\begin{align*}
 \sup_{\gammatwo(T_W)\leq1}|\Tr(W\sigma_\mu)|
 &=1+\bigl(|f_\mu-1/d|-c\bigr)_+\\
 &=1+2(1/d-\mu)_+.
\end{align*}
The absolute value causes no extra branch because
$-1\leq f_\mu\leq1$.  Lemma \ref{lem:duality} proves the corresponding
upper bound for $\alpha$.

For the reverse bound, if $\mu\geq1/d$, the one-dimensional trace testers
$X\mapsto\Tr X$ give output norm $\Tr\sigma_\mu=1$.  If $\mu<1/d$, the
realignment tester gives
\[
 \|(\mathcal R\otimes\mathcal R)(\sigma_\mu)\|_1
 =\frac{d+2}{d}-2\mu
 =1+2(1/d-\mu),
\]
by the direct calculation in \cite[Section 9]{JLN}.  Thus equality holds in
all cases, and realignment attains the nontrivial branch.
\end{proof}

\begin{corollary}[Counterexamples to completeness]\label{cor:counterexample}
For every $d\geq3$ and
\[
 \frac1d\leq\mu<\frac12,
\]
$\sigma_\mu$ is entangled, but no pair of entanglement testers detects it.
In particular, on $\mathbb C^3\otimes\mathbb C^3$,
\[
 \rho_*:=\sigma_{1/3}=\frac{5I-3F}{36}
\]
is an explicit counterexample to the source question.
\end{corollary}

\begin{proof}
For positive $P,Q$, $\Tr(F(P\otimes Q))=\Tr(PQ)\geq0$.  Hence every
separable state has nonnegative flip expectation.  Equation \eqref{eq:f}
is negative when $\mu<1/2$, so $\sigma_\mu$ is entangled.  On the other hand,
Theorem \ref{thm:envelope} gives $\alpha(\sigma_\mu)=1$ when
$\mu\geq1/d$.  Since the tester criterion requires a strict value greater
than one, no tester pair detects these states.

For $d=3$, $F$ has eigenvalue $+1$ with multiplicity six and $-1$ with
multiplicity three.  Thus $\rho_*$ has eigenvalues $1/18$ (six times) and
$2/9$ (three times), has trace one, and
$\Tr(F\rho_*)=-1/3$.  This verifies all density-state and entanglement
claims directly.
\end{proof}

\section*{Verification Notes}

The proof uses the complex $S_1$ unit ball in Lemma \ref{lem:gamma}; its
rank-one extreme points are essential for the exact norm formulas
\eqref{eq:E}--\eqref{eq:G}.  The diagonal lower bound is normalized by $d$,
not $\sqrt d$: a factorization of a $d\times d$ matrix has at most $d$ row
vectors and $d$ column vectors, so the product of their Frobenius norms is at
most $d$ times the factorization constant.

The supplementary script verifies positivity, trace, and flip expectation
for $\rho_*$; scans the invariant support problem for $d=2,\ldots,10$; and
checks the optimizing factorization parameters on 7000 random coefficient
pairs.  It is not used as proof.

\section*{Novelty Check and Scope}

A bounded local, arXiv, and web search on August 11, 2026 used the source id
and title together with ``entanglement testers'', ``Werner'', ``complete'',
``factorization through $\ell_2$'', and the exact open-question wording.  It
found no prior negative answer or exact all-tester Werner envelope.  A 2021
workshop report still states the problem as open.  The six works indexed by
OpenAlex as citing the source through 2025 were checked at title/abstract or
full-text citation context; accessible arXiv sources 2202.13993, 2402.09972,
and 2503.21177 use the paper for adjacent tensor-norm or SIC/realignment
results, not this completeness question.  Novelty is plausible, not
certified.

One bipartite counterexample answers the source's multipartite universal
question negatively.  The result does not address whether testers are
complete for $2\times2$ states or after the source's permitted index
permutation; the latter is already positive in \cite{JLN}.

\section*{Human Review Recommendation}

First verify Lemma \ref{lem:duality}'s normalization and the nuclear-norm
lower bound in Lemma \ref{lem:gamma}.  Then check the one-line support
optimization in Theorem \ref{thm:envelope}.  If those pass, the explicit
$3\times3$ state is a complete negative answer and the full Werner-state
optimality statement is an additional sharp result.

\begin{thebibliography}{9}

\bibitem{JLN}
M. A. Jivulescu, C. Lancien, and I. Nechita,
\emph{Multipartite entanglement detection via projective tensor norms},
Ann. Henri Poincar\'e \textbf{23} (2022), 3791--3838;
arXiv:2010.06365.

\bibitem{Werner}
R. F. Werner,
\emph{Quantum states with Einstein--Podolsky--Rosen correlations admitting a
hidden-variable model},
Phys. Rev. A \textbf{40} (1989), 4277--4281.

\end{thebibliography}

\end{document}
